% !PRINTED 39
\weberpage{39}
Wenn aber $f(x_0)$ verschwindet, so erhält die Entwickelung von $R(x)$ die Form:
\weq{6}{R(x) = A_m(x-x_0)^{m-1/2} + A_{m+1}(x-x_0)^{m-1/2+1} + A_{m+2}(x-x_0)^{m-1/2+2} + \cdots,}
und wenn $f(x)$ vom dritten Grade ist, so gilt in der Umgebung des unendlich fernen Punktes die Entwickelung
\weq{7}{R(x) = C_m x^{-m-3/2} + C_{m+1} x^{-m-5/2} + C_{m+2} x^{-m-7/2} + \cdots}
Aus den Entwickelungen (4) bis (7) können wir die Entwickelungen von $\Omega$ ableiten, wobei eine additive Konstante unbestimmt bleibt.

Die Entwickelung von $\Omega$ ist dann ebenfalls eine Potenzreihe, wozu, wenn $m$ negativ ist, in den Fällen (4) oder (5) noch ein Glied $A_{-1}\log(x-x_0)$, $C_{-1}\log x$ tritt.

\gesperrt{Wir nennen nun $x_0$ einen Punkt erster Gattung von} $\Omega$, wenn $m$ in diesen Entwickelungen nicht negativ ist.

\gesperrt{Dann ist $\Omega$ im Punkte $x_0$ endlich (auch im Falle $x_0 = \infty$).}

\gesperrt{Der Punkt $x_0$ heißt ein Punkt zweiter Gattung von $\Omega$,} wenn in den Entwickelungen (4), (5) $m$ negativ und $A_{-1} = 0$, $C_{-1} = 0$ ist, und in dem Falle der Entwickelung (6), (7) mit negativem $m$.

\gesperrt{Endlich heißt $x_0$ ein Punkt dritter Gattung,} wenn in den Entwickelungen (4) oder (5) $A_{-1}$ oder $C_{-1}$ von Null verschieden ist, wenn also in der Entwickelung von $\Omega$ ein logarithmisches Glied vorkommt.

Die Nullpunkte von $f(x)$ können bei dem Integral (2) nicht von der dritten Gattung sein; sie können aber von der dritten Gattung werden, wenn wir das elliptische Integral in der allgemeinen Form $\int \Phi(x, \sqrt{f(x)})\,dx$ betrachten.

Das Integral $\Omega$ heißt ein \gesperrt{Integral erster Gattung}, wenn es nur Punkte erster Gattung hat. Ein solches Integral hat dann für alle endlichen und unendlichen Werte von $x$ einen endlichen Wert.

$\Omega$ heißt ein \gesperrt{Integral zweiter Gattung}, wenn es nur Punkte erster und zweiter Gattung hat, und

$\Omega$ heißt ein \gesperrt{Integral dritter Gattung}, wenn es neben Punkten erster und zweiter Gattung auch Punkte dritter Gattung hat.
